% mazzi: 07
\chapter{Il tessuto muscolare}

% ---------------------------------------------------------------------------
\section{Generalità}

\textbf{Tessuti}: gruppi di cellule che cooperano con simili caratteristiche morfologiche
$\rightarrow$ \textit{divisione del lavoro}. Quattro tipi: epiteliale, connettivo, muscolare,
nervoso.

\textbf{Muscolo}: organo che genera forza e movimento (es.\ bicipite brachiale, rigonfio se
\textit{attivo}, rilasciato se \textit{inattivo}).
\begin{itemize}
  \item costituito da \textbf{fibre muscolari} in fasci, avvolte da connettivo, raggiunte da vasi e
    fibre nervose motorie;
  \item al ME (sezione trasversale): intreccio di fibre e connettivo; al MO (sezione longitudinale):
    andamento parallelo e nuclei allineati.
\end{itemize}

% ---------------------------------------------------------------------------
\section{I tre tipi di tessuto muscolare}

\subsection{Scheletrico}
\begin{itemize}
  \item cellule \textbf{lunghe, cilindriche, striate, multinucleate}; nuclei periferici (sotto la
    membrana);
  \item nei muscoli scheletrici (es.\ arti), in rapporto con connettivo e tessuto nervoso;
  \item filamenti di actina e miosina $\rightarrow$ \textbf{striature} visibili al MO;
  \item funzioni: muovere/stabilizzare lo scheletro; controllo di ingresso/uscita di digerente,
    respiratorio, urinario; generare calore; proteggere i visceri.
\end{itemize}

\subsection{Cardiaco}
\begin{itemize}
  \item cellule (\textbf{cardiomiociti}) più piccole, ramificate (estremità a Y), con un solo
    \textbf{nucleo} centrale;
  \item interconnesse dai \textbf{dischi intercalari} (trasmissione meccanica ed elettrica);
  \item striate (actina e miosina);
  \item sede: solo il \textbf{cuore} (miocardio) $\rightarrow$ spinge il sangue in circolo, regola
    la pressione ematica.
\end{itemize}

\subsection{Liscio}
\begin{itemize}
  \item cellule \textbf{piccole, affusolate, non striate}, un solo \textbf{nucleo} centrale, senza
    diramazioni;
  \item circondano vasi sanguigni e apparati digerente, respiratorio, urinario, genitale;
  \item funzioni: propulsione di cibo/urina/secrezioni; regolazione del calibro di vie respiratorie
    e vasi $\rightarrow$ flusso ematico tissutale.
\end{itemize}

% ---------------------------------------------------------------------------
\section{Le fibre muscolari lente}

Popolazione di fibre \textbf{lente} del muscolo scheletrico:
\begin{itemize}
  \item \textbf{metabolismo aerobico};
  \item elevato contenuto di \textbf{mioglobina};
  \item numerosi mitocondri;
  \item elevata vascolarizzazione;
  \item \textbf{contrazione prolungata};
  \item colorazione rossastra $\rightarrow$ ``carni rosse''.
\end{itemize}
Adatte all'attività prolungata a bassa intensità (le fibre veloci: contrazioni rapide e potenti ma
poco resistenti alla fatica).

% ---------------------------------------------------------------------------
\section{Le fasce corporee}

\textbf{Fasce}: formazioni di connettivo che danno forza e stabilità, mantengono gli organi in
posizione e distribuiscono vasi, linfatici e nervi. Dalla cute verso l'interno, tre livelli:
\begin{itemize}
  \item \textbf{superficiale}: tra cute e organi sottostanti; tessuto areolare e adiposo (strato
    sottocutaneo = \textbf{ipoderma});
  \item \textbf{profonda}: intelaiatura fibrosa interna di connettivo denso; legata a capsule,
    tendini, legamenti;
  \item \textbf{sottosierosa}: tra membrane sierose e fascia profonda; tessuto areolare.
\end{itemize}
Più in profondità: la \textbf{membrana sierosa} riveste la cavità corporea; a livello della parete
la fascia profonda avvolge le \textbf{coste}; la \textbf{membrana cutanea} è il rivestimento più
esterno.

% ---------------------------------------------------------------------------
\section{Fisiologia della contrazione muscolare}

Sequenza di attivazione della fibra: \textbf{stimolo chimico/elettrico} (via giunzione
neuromuscolare = placca motrice, tra assone del motoneurone e fibra) $\rightarrow$ \textbf{risposta
di membrana} = \textbf{potenziale d'azione} $\rightarrow$ \textbf{risposta intracellulare} =
\textbf{contrazione}.

Registrabile con elettrodi di superficie su bicipite e tricipite (amplificatore biologico + messa a
terra) durante una contrazione volontaria $\rightarrow$ tracciato EMG grezzo e il suo inviluppo
integrato (\textit{integrated EMG}).

\subsection{Elettromiografia (EMG)}
\textbf{EMG}: metodica neurofisiologica per studiare il sistema nervoso periferico dal punto di
vista funzionale. Con elettrodo ad ago nel ventre muscolare (es.\ gamba, avambraccio) si registra
l'attività elettrica delle unità motorie durante la contrazione.

% ---------------------------------------------------------------------------
\section{Tecniche di imaging del tessuto muscolare}

\subsection{Risonanza magnetica (RM)}
Visualizza in dettaglio i singoli muscoli e la loro disposizione.
\begin{itemize}
  \item sezione assiale T1 della coscia: vasto laterale (1), vasto mediale e intermedio (2), retto
    femorale (3), sartorio (4), adduttore lungo (5), adduttore breve (6), gracile (7), semitendinoso
    (8), grande gluteo (9);
  \item sezione coronale T1: bicipite femorale (1), vasto laterale (2), gemello inferiore (3),
    otturatorio interno (4), gruppo adduttore (5), gracile (6).
\end{itemize}

\subsection{Spettroscopia a risonanza magnetica (MRS)}
\textbf{\textsuperscript{31}P MRS}: studia il metabolismo energetico del muscolo a riposo
(\textit{at rest}). Spettro al polpaccio con picchi di \textbf{fosfato inorganico (Pi)},
\textbf{fosfodiesteri (PDE)}, \textbf{fosfocreatina (PCr)} (il più intenso) e i tre fosfati
dell'\textbf{ATP} ($\gamma$, $\alpha$, $\beta$). Acquisizione guidata da un'immagine RM trasversale
del polpaccio (1,5 T) con un riquadro di interesse $\rightarrow$ dati di spettroscopia protonica
(\textsuperscript{1}H MRS) sui lipidi muscolari.

% ---------------------------------------------------------------------------
\section{Biopsia muscolare}

\textbf{Biopsia muscolare}: valuta le anomalie muscolo-scheletriche e diagnostica le malattie del
tessuto muscolare. Prelievo con dispositivo bioptico (es.\ strumento a scatto ``Magnum'') o per via
chirurgica (piccolo campo operatorio con divaricatori).

\rubrica{Rischi}
\begin{itemize}
  \item lividi e dolore al sito;
  \item sanguinamento prolungato;
  \item infezione del sito.
\end{itemize}

\rubrica{Esame microscopico}
Al MO: morfologia longitudinale (fibre striate e parallele) e sezione trasversale (profili
poligonali delle fibre, delimitati dal \textbf{perimisio} = connettivo dei fascicoli, attraversati
da un \textbf{capillare} nello spazio interfascicolare).

% ---------------------------------------------------------------------------
\section{Alterazioni patologiche}

L'esame istologico può rivelare:
\begin{itemize}
  \item \textbf{perdita di striature}: disorganizzazione dell'apparato contrattile;
  \item \textbf{degenerazione con infiltrati infiammatori}: processo patologico in atto sulle fibre.
\end{itemize}
